\section{Finite fields, quadratic reciprocity, and sums of squares}
\label{sec:quadratic}

This section develops three connected parts of number theory: the elementary
structure of the nonzero residues modulo a prime, quadratic reciprocity, and
representations by two and four squares.  The conventional topics and proof
traditions may be compared with \cite{lynn-qr,kumar-qr,conrad-descent}.
The proofs below include their own constructions and bounds.  In particular,
``descent'' always means an explicitly bounded iteration, not an appeal to an
unbounded least-number principle.

\subsection{Residue arithmetic and polynomial roots}

For a prime $p$, write $\mathbf F_p$ as notation for the integers
$0,\ldots,p-1$ with addition and multiplication followed by remainder modulo
$p$. This is a finite representation, not a new sort of the arithmetic.
An inverse of $a\ne0$ is the first $b<p$ with $ab\equiv1\pmod p$;
B\'ezout's identity proves that this bounded search succeeds.  In particular,
$ab=0$ in $\mathbf F_p$ implies $a=0$ or $b=0$.  Polynomials are finite
coefficient arrays and their values are computed by Horner's recurrence.
Trailing zero coefficients are removed by a bounded scan.  All assertions
about roots mean finite tests over $0,\ldots,p-1$.

\begin{theorem}[Polynomial root bound]\label{qr:roots}
A nonzero polynomial of degree $d$ over $\mathbf F_p$ has at most $d$
distinct roots.
\end{theorem}
\begin{proof}
For a constant nonzero polynomial there are no roots.  If a polynomial $f$
of positive degree has no root, the conclusion also holds.  Otherwise choose
the first root $a<p$.  Synthetic division computes an array $g$ of degree
$d-1$ such that
\[
 f(X)=(X-a)g(X)+f(a)=(X-a)g(X).
\]
For completeness, if $f=\sum_{i=0}^d c_iX^i$, set $b_{d-1}=c_d$ and,
working downwards, $b_{i-1}=c_i+ab_i$ for $1\le i<d$.  The coefficient
of $X^i$ in $(X-a)\sum_{j=0}^{d-1}b_jX^j$ is $b_{i-1}-ab_i=c_i$;
the remaining constant term agrees because evaluation at $a$ gives zero.
If $b\ne a$ is another root of $f$, cancellation in $\mathbf F_p$ gives
$g(b)=0$.  Thus the roots of $f$ consist of $a$ and at most $d-1$ further
roots, by induction.

Here the induction can be made quantifier-free without a hidden polynomial
sort. Fix $p$ and use the characteristic function which checks the assertion
for every coefficient array of length at most $d+1$ with entries below $p$.
Such arrays, padded with zeros, have codes below $p^{d+1}$.  The induction
step just given reduces the relevant array to one of length at most $d$;
all root counts are sums of $p$ zero-one values.  This is bounded induction
compiled into the quantifier-free induction scheme of \ipra.
\end{proof}

For $1\le a<p$, define $\operatorname{ord}_p(a)$ to be the first positive
$e\le p-1$ with $a^e\equiv1\pmod p$. Fermat's little theorem proves that
the search succeeds; exponentiation is always the ordinary primitive
recursive operation, or its remainder-reduced version.

\begin{lemma}[Orders and coprime products]\label{qr:orders}
If $d=\operatorname{ord}_p(a)$, then
\[
 a^k\equiv1\pmod p\quad\Longleftrightarrow\quad d\mid k
 \qquad(k\in\N),
\]
and $d\mid p-1$.  If $u,v$ have coprime orders $s,t$, respectively,
then $uv$ has order $st$.
\end{lemma}
\begin{proof}
Divide $k$ by $d$, writing $k=qd+r$ with $0\le r<d$.  If $a^k=1$,
then $a^r=1$, so the minimality in the bounded definition of $d$ forces
$r=0$.  The reverse implication follows by exponentiation. Apply this to
$k=p-1$ for the second assertion.  For the last assertion, $(uv)^{st}=1$.
If $(uv)^k=1$, raising to $t$ gives $u^{kt}=1$, hence $s\mid kt$.
Coprimality gives $s\mid k$ by B\'ezout. Similarly $t\mid k$, and therefore
$st\mid k$. The least positive such exponent is $st$.
\end{proof}

\begin{theorem}[Primitive roots modulo a prime]\label{qr:primitive}
For every prime $p$ there is $g<p$ whose powers run through all the
nonzero residues modulo $p$. A primitive recursive function can choose
such a $g$ by searching $1\le g<p$.
\end{theorem}
\begin{proof}
For $p=2$, take $g=1$. Put $N=p-1$ and factor $N$ as
$\prod_{i<r}q_i^{e_i}$ with distinct primes $q_i$. For each $i$, the
polynomial $X^{N/q_i}-1$ has degree $N/q_i<N$. By the root bound it
cannot vanish at all $N$ nonzero residues. Choose the first $a_i<p$
for which $a_i^{N/q_i}\ne1$, and set
\[
 b_i=a_i^{N/q_i^{e_i}}\pmod p.
\]
The order of $b_i$ divides $q_i^{e_i}$. If it were a proper divisor,
it would divide $q_i^{e_i-1}$, contradicting
$b_i^{q_i^{e_i-1}}=a_i^{N/q_i}\ne1$. Thus its order is exactly
$q_i^{e_i}$. The preceding lemma shows that $g=\prod_{i<r}b_i$
has order $N$. If $g^i=g^j$ for $0\le i<j<N$, cancellation gives
$g^{j-i}=1$, contradicting the order. Consequently these $N$ powers
are distinct, and finite counting shows that they exhaust the $N$
nonzero residues.

The factor list has length at most $N$, every residue search has $p$
candidates, and each exponent is at most $N$. Alternatively the final
function may simply test every $g<p$ for the $N$ required inequalities;
the argument proves that this single bounded search succeeds.
\end{proof}

\begin{corollary}[Orders of powers and the number of primitive roots]
\label{qr:primitive-count}
For a primitive root $g$ and $N=p-1$,
\[
 \operatorname{ord}_p(g^k)=\frac{N}{\gcd(N,k)}.
\]
There are exactly $\varphi(N)$ primitive roots modulo $p$. More generally,
if $b\ne0$, $k\ge1$, $b=g^t$ with $0\le t<N$, and $d=\gcd(k,N)$,
then $x^k=b$ has no nonzero solution when $d\nmid t$, and exactly $d$
nonzero solutions when $d\mid t$.
\end{corollary}
\begin{proof}
The equation $(g^k)^j=1$ is equivalent to $N\mid kj$. Dividing by
$d=\gcd(N,k)$ and using coprimality shows that this is equivalent to
$N/d\mid j$. This proves the order formula. Since $g^0,\ldots,g^{N-1}$
list the units once, exactly those exponents coprime to $N$ give order
$N$, proving the count. Finally, write $x=g^j$; then $x^k=b$ is
$kj\equiv t\pmod N$. Necessarily $d\mid t$. When $d\mid t$, division
by $d$ gives a congruence with invertible coefficient modulo $N/d$,
so it has one class $j_0$ there. Its $d$ distinct representatives modulo
$N$ are $j_0+\ell N/d$, $0\le\ell<d$.
\end{proof}

\subsection{Quadratic characters and reciprocity}

For an odd prime $p$ and an integer $a$, define the Legendre symbol by
the finite decision
\[
 \chi_p(a)=\left(\frac ap\right)=
 \begin{cases}
 0,&a\equiv0\pmod p,\\
 1,&a\not\equiv0\pmod p\ \text{and some }1\le x<p\text{ satisfies }x^2\equiv a,\\
 -1,&\text{otherwise}.
 \end{cases}
\]
The word ``otherwise'' is justified by a complete finite search, not by
excluded middle for an arbitrary proposition. The output is encoded as a
signed integer. Put $h=(p-1)/2$.

\begin{lemma}[The square residues]\label{qr:squares-half}
The $h$ residues $1^2,2^2,\ldots,h^2$ are distinct, and they are
exactly the nonzero square residues. Each has precisely two square roots
modulo $p$.
\end{lemma}
\begin{proof}
If $i^2=j^2$, then $(i-j)(i+j)=0$. For $1\le i,j\le h$, the second
factor is a positive integer less than $p$, so it cannot vanish modulo
$p$; hence $i=j$. Every nonzero residue $x$ has either $x\le h$ or
$p-x\le h$, and $x^2=(p-x)^2$. This proves exhaustion. The same
factorization shows that the only roots of $X^2-i^2$ are $i$ and $-i$;
they differ because $p$ is odd and $i\ne0$.
\end{proof}

\begin{theorem}[Euler's criterion]\label{qr:euler-criterion}
For an odd prime $p$ and $p\nmid a$,
\[
 a^{(p-1)/2}\equiv\chi_p(a)\pmod p.
\]
\end{theorem}
\begin{proof}
If $a=b^2\ne0$, Fermat's little theorem gives $a^h=b^{p-1}=1$.
Thus all $h$ nonzero square residues are roots of $X^h-1$. The root
bound permits no further roots, so a nonsquare $a$ satisfies $a^h\ne1$.
On the other hand, $(a^h)^2=a^{p-1}=1$. The identity
$(a^h-1)(a^h+1)=0$ forces $a^h=-1$. These are exactly the two
nonzero values specified by $\chi_p$.
\end{proof}

\begin{corollary}[Multiplicativity]\label{qr:character}
For all integers $a,b$,
\[
 \chi_p(ab)=\chi_p(a)\chi_p(b).
\]
\end{corollary}
\begin{proof}
If either residue is zero, both sides are zero. Otherwise Euler's criterion
makes the two sides congruent modulo $p$. Both are in $\{-1,1\}$.
Since an odd prime does not divide $2$, these two values are distinct
modulo $p$, and congruence implies equality as signed integers.
\end{proof}

\begin{theorem}[Gauss's lemma]\label{qr:gauss}
Let $p\nmid a$. For $1\le i\le h$, let $r_i$ be the residue of
$ai$ in $\{1,\ldots,p-1\}$, and let
\[
 m=\#\{i:1\le i\le h,\ 2r_i>p\}.
\]
Then $\chi_p(a)=(-1)^m$.
\end{theorem}
\begin{proof}
Define $s_i=r_i$ when $r_i\le h$ and $s_i=p-r_i$ otherwise.
These numbers lie in $\{1,\ldots,h\}$. If $s_i=s_j$, then
$ai\equiv aj$ or $ai\equiv-aj$. Cancellation gives $i\equiv j$
or $i\equiv-j$. The first gives $i=j$; the second is impossible
because $2\le i+j\le p-1$. Thus $s_1,\ldots,s_h$ is a permutation
of $1,\ldots,h$. Multiplying all congruences gives
\[
 a^h h!\equiv\prod_{i=1}^h r_i
       \equiv(-1)^m\prod_{i=1}^h s_i=(-1)^m h!\pmod p.
\]
No factor of $h!$ is divisible by $p$, so cancel it and use Euler's
criterion. The comparison of the two signs is justified as in the previous
corollary. Every permutation assertion here is a bounded statement about
$h$ entries; the product equality follows from counting each possible
entry once.
\end{proof}

\begin{lemma}[The floor-sum form]\label{qr:eisenstein}
If $a$ is a positive odd integer coprime to the odd prime $p$, then
\[
 \chi_p(a)=(-1)^{\displaystyle\sum_{i=1}^{h}\lfloor ai/p\rfloor}.
\]
\end{lemma}
\begin{proof}
Use the residues $r_i$ and folded residues $s_i$ of Gauss's lemma and
put $K=\sum_i\lfloor ai/p\rfloor$. Summing $ai=p\lfloor ai/p\rfloor+r_i$
and replacing each negative-side residue by $p-s_i$ yields the exact
integer identity
\[
 (a-1)\sum_{i=1}^{h}i
   =p(K+m)-2\sum_{\{i:r_i>h\}}s_i.
\]
Here we used $\sum_i s_i=\sum_i i$, already proved by the finite
permutation. Since $a-1$ is even and $p$ is odd, $K+m$ is even.
Therefore $(-1)^K=(-1)^m$, and Gauss's lemma proves the assertion.
The identity, rather than any geometric parity intuition, is the
quantifier-free invariant needed for formalization.
\end{proof}

\begin{theorem}[Quadratic reciprocity]\label{qr:reciprocity}
For distinct odd primes $p,q$,
\[
 \left(\frac qp\right)\left(\frac pq\right)
       =(-1)^{((p-1)/2)((q-1)/2)}.
\]
\end{theorem}
\begin{proof}
Put $h=(p-1)/2$, $k=(q-1)/2$. Consider all $hk$ pairs
$(i,j)$ with $1\le i\le h$, $1\le j\le k$. None satisfies
$qi=pj$: otherwise $p\mid i$, impossible because $0<i<p$.
For fixed $i$, the number with $pj<qi$ is $\lfloor qi/p\rfloor$.
Indeed $qi/p<q/2=k+1/2$, and $qi/p$ is not integral, so this floor
lies between $0$ and $k$. For fixed $j$, the number with $qi<pj$
is similarly $\lfloor pj/q\rfloor$. The decidable alternatives
$pj<qi$ and $qi<pj$ partition the rectangle. Consequently
\[
 \sum_{i=1}^{h}\left\lfloor\frac{qi}{p}\right\rfloor+
 \sum_{j=1}^{k}\left\lfloor\frac{pj}{q}\right\rfloor=hk.
\]
Apply the floor-sum formula twice and multiply the signs. The counting
identity is proved arithmetically by two nested bounded sums, so no
geometric theorem or principle concerning an infinite set has entered.
\end{proof}

\begin{theorem}[The two supplementary laws]\label{qr:supplements}
For every odd prime $p$,
\[
 \left(\frac{-1}{p}\right)=(-1)^{(p-1)/2},
 \qquad
 \left(\frac2p\right)=(-1)^{(p^2-1)/8}.
\]
In particular, $-1$ is a square precisely when $p\equiv1\pmod4$,
and $2$ is a square precisely when $p\equiv1$ or $7\pmod8$.
\end{theorem}
\begin{proof}
The first identity follows immediately by applying Euler's criterion to
$-1$. For the second, the residues of $2,4,\ldots,2h$ are those
integers themselves because $2h=p-1$. In Gauss's lemma, a residue
$2i$ is on the negative side exactly when $4i>p$. Hence
\[
 m=h-\lfloor p/4\rfloor.
\]
Write the odd integer $p$ in one of the four forms $8t+1,8t+3,8t+5,8t+7$.
The corresponding values of $m$ are $2t,2t+1,2t+1,2t+2$.
The integers $(p^2-1)/8$ in those four cases have the same parities:
$8t^2+2t$, $8t^2+6t+1$, $8t^2+10t+3$, and $8t^2+14t+6$.
Thus Gauss's sign is exactly the displayed sign. These four branches
also give the stated residue-class criteria.
\end{proof}

\subsection{Two squares: prime representation and exact classification}

Write $N_2(x,y)=x^2+y^2$ for signed integers $x,y$.  A pair of
nonnegative square roots representing $n$ can always be sought in the
finite box $0\le x,y\le n$.  We use signed coordinates during the
constructions because subtraction makes the identities shorter.

\begin{lemma}[The two-square product identity]\label{qr:two-product}
For all integers $x,y,r,s$,
\[
 (x^2+y^2)(r^2+s^2)=(xr+ys)^2+(xs-yr)^2.
\]
Consequently products of sums of two squares are sums of two squares.
\end{lemma}
\begin{proof}
Expanding the right side gives $x^2r^2+2xyrs+y^2s^2+x^2s^2-2xyrs+y^2r^2$.
The mixed terms cancel and the remaining four terms factor as claimed.
This is a polynomial identity over the signed-integer coding, hence an
identity provable from the elementary ring equations of \ipra.
\end{proof}

\begin{lemma}[A strict, bounded two-square descent]\label{qr:two-descent}
Let $p$ be prime and $1<m<p$. From $x^2+y^2=mp$ one can compute
integers $u,v$ and $1\le m'<m$ with $u^2+v^2=m'p$.
\end{lemma}
\begin{proof}
Choose balanced residues $r\equiv x$ and $s\equiv y\pmod m$ with
$|r|,|s|\le m/2$. They are computed by taking ordinary remainders
and subtracting $m$ exactly when twice the remainder exceeds $m$.
Since $r^2+s^2\equiv x^2+y^2\equiv0\pmod m$, write
$r^2+s^2=mm'$. The estimate $r^2+s^2\le m^2/2$ gives
$m'\le m/2<m$. If $m'=0$, both $r$ and $s$ vanish, so $m$ divides
$x$ and $y$, hence $m^2\mid mp$. Cancellation yields $m\mid p$,
contrary to $1<m<p$. Thus $m'>0$.

The numerators $xr+ys$ and $xs-yr$ are divisible by $m$, by reducing
them modulo $m$. Define
\[
 u=(xr+ys)/m,\qquad v=(xs-yr)/m.
\]
The product identity gives $u^2+v^2=m'p$. All divisions displayed here
are exact integer divisions established before they are performed.
A state has $m<p$ and $|x|,|y|<p$; the new state has the same bounds.
The intermediate products have absolute value at most $p^2$, and a
sum of two has absolute value at most $2p^2$.
\end{proof}

\begin{theorem}[Fermat's two-square theorem for primes]\label{qr:two-prime}
A prime $p$ is a sum of two integer squares if and only if $p=2$ or
$p\equiv1\pmod4$. In either positive case a primitive recursive
function returns a pair $(x,y)$ with $0\le x,y\le p$ and
$x^2+y^2=p$.
\end{theorem}
\begin{proof}
The prime $2$ is $1^2+1^2$. Suppose $p\equiv1\pmod4$.
The first supplementary law, together with the defining bounded search
for the Legendre symbol, supplies $a$ with $a^2\equiv-1\pmod p$.
Replacing $a$ by $p-a$ if needed gives $1\le a\le(p-1)/2$.
Thus $a^2+1=mp$ for some $1\le m<p$; the upper bound follows
from $a^2+1<p^2$. If $m>1$, use the preceding descent step.
Repeat this operation exactly $p$ times, freezing the state when $m=1$.
At every active step the positive multiplier drops by at least one;
therefore it must equal $1$ by the end. Absolute values of the final
coordinates give the requested pair.

Conversely, squares are $0$ or $1$ modulo $4$. If an odd prime is
a sum of two squares, exactly one square is odd, and the sum is
$1$ modulo $4$. This covers all primes.

To make the fuel argument explicit, if $m_0<p$ is the initial
multiplier and $m_t$ the multiplier after $t$ frozen-or-descending
steps, the bounded invariant is
$m_t\le\max(1,m_0-t)$, where subtraction is truncated in the
arithmetic coding. Quantifier-free induction proves it, and at $t=p$
it gives $m_p=1$. Thus the proof constructs the solution; it does not
merely negate the possibility of infinite descent.
\end{proof}

\begin{lemma}[Primes congruent to three modulo four]\label{qr:two-obstruction}
If $q$ is prime, $q\equiv3\pmod4$, and $q\mid x^2+y^2$, then
$q\mid x$ and $q\mid y$. If $n=x^2+y^2>0$, its $q$-adic
valuation $v_q(n)$ is even.
\end{lemma}
\begin{proof}
If $q\nmid y$, divide the congruence $x^2+y^2=0$ by $y^2$
in $\mathbf F_q$. It gives a square root of $-1$, contradicting the
first supplementary law. Thus $q\mid y$, and then $q\mid x^2$
forces $q\mid x$ by primality.

For the valuation statement, choose $t$ to be the smaller of the two
valuations of the nonzero coordinates, ignoring a zero coordinate.
At least one coordinate is nonzero because $n>0$. Equivalently, choose
the greatest $t\le n$ for which $q^t$ divides both coordinates;
there cannot be a larger relevant exponent, since a nonzero coordinate
has absolute value at most $n$ and $2^{n+1}>n$. Write
$x=q^t u$, $y=q^t v$. At least one of $u,v$ is not divisible by $q$.
The first part implies $q\nmid u^2+v^2$. Since
$n=q^{2t}(u^2+v^2)$, unique factorization gives $v_q(n)=2t$.
All the valuation selections used here are bounded by $n$.
\end{proof}

\begin{theorem}[The sum-of-two-squares classification]\label{qr:two-classification}
For $n>0$, the following conditions are equivalent:
\begin{enumerate}
\item $n=x^2+y^2$ for integers $x,y$;
\item every prime $q\equiv3\pmod4$ has even exponent in the
prime factorization of $n$.
\end{enumerate}
There is a primitive recursive decision procedure for this condition,
and a primitive recursive witness function returning nonnegative
$x,y\le n$ whenever it holds. Also $0=0^2+0^2$.
\end{theorem}
\begin{proof}
Necessity is the preceding lemma, applied to every prime factor of $n$.
For sufficiency, take the finite prime factorization. Each factor $2$
has representation $(1,1)$, and each prime $p\equiv1\pmod4$ has
a representation from the prime theorem. A prime factor
$q\equiv3\pmod4$ occurs to an even exponent $2e$, so the whole
factor $q^{2e}$ has representation $(q^e,0)$. Multiply these finitely
many representations using the two-square identity, starting with
$(1,0)$. The norm invariant states after each multiplication that the
current pair represents the product of the factors already processed.
The final product is $n$. Taking absolute values gives nonnegative
coordinates, each at most $n$.

The decision is a scan over primes $q\le n$ and their valuations.
For a particularly simple total witness function, scan the entire box
$\{0,\ldots,n\}^2$ and return its first representation, with $(0,0)$
as a failure value when none exists. The proof establishes success
under condition (2); the search itself has exactly $(n+1)^2$ candidates.
Thus the internal statement may use only the equality
$N_2(\operatorname{twoSq}(n))=n$ and a zero-one condition checking (2).
\end{proof}

\subsection{Four squares by finite descent}

Write $N_4(x)=x_0^2+x_1^2+x_2^2+x_3^2$. For two signed integer
quadruples define the bilinear quadruple $B(x,r)$ by
\begin{align*}
 B_0&=x_0r_0+x_1r_1+x_2r_2+x_3r_3,\\
 B_1&=x_0r_1-x_1r_0-x_2r_3+x_3r_2,\\
 B_2&=x_0r_2+x_1r_3-x_2r_0-x_3r_1,\\
 B_3&=x_0r_3-x_1r_2+x_2r_1-x_3r_0.
\end{align*}

\begin{lemma}[Euler's four-square identity]\label{qr:four-product}
For the displayed bilinear map,
\[
 N_4(B(x,r))=N_4(x)N_4(r).
\]
Moreover, if $r_i\equiv x_i\pmod m$ for all $i$ and
$m\mid N_4(x)$, then $m\mid B_i(x,r)$ for every $i$.
\end{lemma}
\begin{proof}
In the sum of the four squared coordinates, every term $x_i^2r_j^2$
occurs once. Every mixed term has a second occurrence with the opposite
sign; cancellation leaves $\sum_{i,j}x_i^2r_j^2=N_4(x)N_4(r)$.
For the divisibility assertion substitute $r_i=x_i$ modulo $m$.
The first coordinate becomes $N_4(x)$, and each remaining coordinate
becomes zero by commutativity of multiplication. This also explains the
particular sign convention in $B$: the divisibility needed for descent
is built directly into the identity.
\end{proof}

\begin{lemma}[A small represented multiple of an odd prime]\label{qr:four-seed}
For every odd prime $p$ one can find $a,b\le(p-1)/2$ and
$1\le m<p$ such that
\[
 a^2+b^2+1^2+0^2=mp.
\]
\end{lemma}
\begin{proof}
Put $h=(p-1)/2$. The residues $a^2$ for $0\le a\le h$ are
pairwise distinct: the usual factorization of $a^2-b^2$ gives this,
including the zero case. The residues $-1-b^2$ for $0\le b\le h$
are also pairwise distinct. Each list has $h+1$ elements, whereas
there are only $p=2h+1$ residues. If they had no common element,
the concatenated list would contain $p+1$ distinct residues in a box
of size $p$, contrary to finite counting. A scan of the
$(h+1)^2$ pairs therefore returns an intersection, so
$p\mid a^2+b^2+1$. This positive number is at most
$2h^2+1=(p^2-2p+3)/2<p^2$ for $p\ge3$. Division by $p$
gives the required multiplier.

The counting contradiction eliminates a decidable proposition: complete
search either finds a pair or certifies all pairs unequal. Hence it
constructs a pair intuitionistically. No infinitary pigeonhole principle
is used.
\end{proof}

\begin{lemma}[A strict four-square descent, including the even case]
\label{qr:four-descent}
If $p$ is an odd prime, $1<m<p$, and $N_4(x)=mp$, one can
compute $1\le m'<m$ and $y$ with $N_4(y)=m'p$.
\end{lemma}
\begin{proof}
First suppose $m$ is even. Since $mp$ is even, an even number of
the four coordinates is odd. There are therefore pairs of equal parity:
by a permutation, arrange $x_0\equiv x_1\pmod2$ and
$x_2\equiv x_3\pmod2$. This can be found by testing the $24$
permutations. Set
\[
 y=\left(\frac{x_0+x_1}{2},\frac{x_0-x_1}{2},
          \frac{x_2+x_3}{2},\frac{x_2-x_3}{2}\right).
\]
The coordinates are integers and
$N_4(y)=N_4(x)/2=(m/2)p$. Thus take $m'=m/2$, which is
positive and smaller than $m$.

Now suppose $m$ is odd. Choose balanced residues
$r_i\equiv x_i\pmod m$ with $|r_i|\le(m-1)/2$.
Their norm is divisible by $m$, so write $N_4(r)=mm'$.
The strict estimate
\[
 N_4(r)\le4\left(\frac{m-1}{2}\right)^2=(m-1)^2<m^2
\]
gives $m'<m$. If $m'=0$, all $r_i$ vanish, so every $x_i$
is divisible by $m$. Then $m^2\mid mp$, hence $m\mid p$,
contrary to $1<m<p$. Thus $m'\ge1$. The previous lemma makes
all coordinates of $B(x,r)$ divisible by $m$. Define
$y_i=B_i(x,r)/m$. Its norm is
\[
 N_4(y)=\frac{N_4(x)N_4(r)}{m^2}
        =\frac{(mp)(mm')}{m^2}=m'p.
\]
This completes both branches.

The even branch is essential for this proof: balanced residues for an
even modulus only give $N_4(r)\le m^2$, which need not be strict.
The parity operation handles that case directly. For every valid state,
$|x_i|<p$ because $N_4(x)<p^2$. The same holds for the output,
and every bilinear numerator has absolute value at most $4p^2$.
These bounds cover the whole finite iteration and all signed arithmetic.
\end{proof}

\begin{theorem}[Lagrange's four-square theorem]\label{qr:four-squares}
Every $n\in\N$ is a sum of four squares of nonnegative integers.
There is a primitive recursive function $\operatorname{fourSq}$ such
that, provably in \ipra,
\[
 \sum_{i=0}^3\operatorname{fourSq}_i(n)^2=n,
 \qquad \operatorname{fourSq}_i(n)\le n.
\]
\end{theorem}
\begin{proof}
For an odd prime $p$, the seed lemma supplies a represented multiple
$mp$ with $1\le m<p$. Apply the descent step when $m>1$ and
freeze the state when $m=1$. Exactly $p$ iterations suffice, by the
same invariant $m_t\le\max(1,m_0-t)$ used in the two-square proof.
This yields a representation of $p$. The prime $2$ is
$1^2+1^2+0^2+0^2$.

The numbers $0$ and $1$ have the evident representations. For $n>1$,
compute its prime factor list and multiply the prime representations
using $B$, starting from $(1,0,0,0)$. The four-square identity proves,
by induction over this list, that the current norm is the product of
the already processed factors. At the end the norm is $n$; take
absolute values of the coordinates. A nonnegative coordinate with
square at most $n$ is at most $n$.

All fuel bounds are explicit: factorization uses at most $n$ extracted
factors, each prime routine uses at most $p\le n$ descent steps,
and the final product uses at most $n$ bilinear operations. A simpler
specification of the witness function is the first quadruple in
$\{0,\ldots,n\}^4$ of norm $n$, with a default zero quadruple.
This is a bounded scan of $(n+1)^4$ candidates; the constructive proof
above establishes its success for every input. The universal theorem
therefore has a quantifier-free matrix in the language with this
primitive recursive witness function.
\end{proof}

\begin{corollary}[Four squares over the nonnegative rationals]
\label{qr:rational-four}
For integers $a\ge0$ and $b\ge1$, the rational $a/b$ is a sum of
four rational squares, with a primitive recursive construction from
$(a,b)$.
\end{corollary}
\begin{proof}
Represent $ab=x_0^2+x_1^2+x_2^2+x_3^2$ by the preceding theorem.
Then $a/b=\sum_i(x_i/b)^2$. This statement is an exact equality
after multiplication by $b^2$, so it needs no analytic theory of
rational or real numbers. The numerators satisfy $x_i\le ab$ and
the common positive denominator is the input $b$.
\end{proof}

\begin{proposition}[The three-square obstruction and optimality]
\label{qr:three-obstruction}
No integer of the form $4^a(8b+7)$, with $a,b\in\N$, is a sum
of three integer squares. In particular four squares are sometimes
necessary: $7$ cannot be expressed using three squares.
\end{proposition}
\begin{proof}
The square residues modulo $8$ are $0,1,4$. A sum of three such
residues is never $7$: with no $4$ it is at most $3$; with one $4$
it is $4,5$, or $6$; with two $4$ it is $0$ or $1$ modulo $8$;
with three it is $4$. Thus $8b+7$ is not a sum of three squares.
If a sum of three squares is divisible by $4$, all three coordinates
are even, since a square is $0$ or $1$ modulo $4$ and a sum of at
most three ones is zero modulo $4$ only when there are none. Divide
all three coordinates by $2$ and the equation by $4$, and repeat
this operation exactly $a$ times. A representation of
$4^a(8b+7)$ would thereby yield the already impossible representation
of $8b+7$. This is a bounded descent on the supplied exponent $a$.
No converse to this obstruction, and thus no proof here of the full
three-square theorem, is asserted.
\end{proof}

\subsection{Internal arithmetic statements and Lean obligations}
\label{qr:formalization}

All functions in this section are given by composition, primitive recursion,
or explicitly bounded search. For clarity, a common implementation can use
these numerical interfaces:
\begin{itemize}
\item $\operatorname{eval}(p,c,d,x)$: Horner evaluation of a base-$p$
coefficient code $c<p^{d+1}$; $\operatorname{rootCount}(p,c,d)$ is the
sum of its zero tests over $x<p$.
\item $\operatorname{order}(p,a)$: the first $1\le e\le p-1$ with
$a^e\bmod p=1$, defaulting to $0$; $\operatorname{generator}(p)$:
the first $1\le g<p$ with order $p-1$, defaulting to $0$.
\item $\operatorname{legendre}(p,a)$: the signed three-valued finite
search above; $\operatorname{fold}(p,r)=\min(r,p-r)$ for $r<p$;
$\operatorname{negCount}(p,a)$: the sum of the comparisons
$2(ai\bmod p)>p$ for $1\le i\le(p-1)/2$.
\item $\operatorname{balance}(m,x)$: the signed balanced remainder;
$\operatorname{twoStep}$ and $\operatorname{fourStep}$: the exact
formulas in the descent lemmas, with invalid states sent to a fixed
zero state. Valid states have $1\le m<p$ and norm $mp$.
\item $\operatorname{twoSq}(n)$ and $\operatorname{fourSq}(n)$:
the bounded lexicographic searches specified in the main theorems.
A tuple is coded in base $n+1$, so its search bound is respectively
$(n+1)^2$ and $(n+1)^4$.
\end{itemize}
The state records and coefficient codes are numerical encodings.  In
particular, a statement about a list means a statement about the bounded
digit tests of its code. The arithmetic predicates ``prime'', ``coprime'',
``has order'', and ``all $3\bmod4$ exponents are even'' are zero-one
primitive recursive functions. Signed equations can equally be converted
to equations between natural numbers by moving negative summands to the
opposite side.

A Lean formalization should have two distinct layers.  The first is a
semantic development on \texttt{Nat}, coded signed integers, and finite
arrays; the second constructs or checks the corresponding derivations
in the explicit syntax of \ipra. Proving a result in Lean's ambient
logic, even without invoking \texttt{Classical}, is not by itself the
second layer. The specific proof obligations are as follows.
\begin{enumerate}
\item Prove the Horner synthetic-division identity, root-set inclusion,
and counting inequality. Instantiate bounded induction on degree with
all coefficient codes below $p^{d+1}$.
\item Prove the remainder/exponent identities for orders, the finite
root-count witness for each factor of $p-1$, and the coprime-order product
lemma. No general finite-field classification theorem is needed.
\item Prove that folded residues give a bounded permutation. Derive
both the product identity for Gauss's lemma and the exact sum identity
for floor parity. Prove the reciprocity rectangle identity as equality
of two finite sums of zero-one comparisons.
\item Normalize the polynomial identities $N_2$ and $N_4$ in a
commutative-ring semantics, then supply equation-level certificates
for their use in \ipra. Prove exact divisibility before constructing
any quotient in a descent step.
\item For four-square descent, formalize the equal-parity pairing of
four integers separately from the odd-modulus balanced-residue estimate.
This prevents a weak inequality from being mistaken for strict descent.
\item Compile the frozen-loop invariant
$m_t\le\max(1,m_0-t)$ and the norm invariant to quantifier-free
induction. Finish with bounded-search success and its primitive
recursive definition certificate.
\end{enumerate}
This list is a completion route, not a claim that the Lean implementation
of all these obligations has already been completed. It is deliberately
compatible with reusing the repository's folded-residue, Gauss-product,
reciprocity, and supplementary-law derivations once their exact interfaces
and assumptions have been checked.
